\pagestyle{empty}
\thispagestyle{empty}
\setlength{\parindent}{0pt}
\setlength{\tabcolsep}{3pt}
\setlist[itemize]{nosep,leftmargin=6mm}
\setlist[enumerate]{nosep,leftmargin=6.5mm}

\providecommand{\wk}{\cellcolor{black!15}}
\providecommand{\blankwk}{}

\begin{center}
МИНИСТЕРСТВО НАУКИ И ВЫСШЕГО ОБРАЗОВАНИЯ\\
РОССИЙСКОЙ ФЕДЕРАЦИИ

\vspace{6pt}

Федеральное государственное автономное образовательное учреждение высшего образования\\
«МОСКОВСКИЙ ПОЛИТЕХНИЧЕСКИЙ УНИВЕРСИТЕТ»\\
(МОСКОВСКИЙ ПОЛИТЕХ)

\vspace{12pt}

\textbf{ЗАДАНИЕ НА ВЫПУСКНУЮ КВАЛИФИКАЦИОННУЮ РАБОТУ}

\vspace{6pt}

по направлению 09.03.03 Прикладная информатика\\
Образовательная программа (профиль) «Корпоративные информационные системы»
\end{center}

\vspace{10pt}

\textbf{ТЕМА ВКР}

«Разработка информационной системы для создания туристических маршрутов с геопривязкой фотографий и интерактивным обменом на основе микросервисной архитектуры»

\vspace{8pt}

\textbf{ПРАКТИЧЕСКИЙ РЕЗУЛЬТАТ}

Программная система Guide Helper для создания, публикации и сопровождения туристических маршрутов с геопривязанными фотографиями, интерактивной картой и средствами совместного доступа.

\vspace{4pt}

\textbf{Назначение}

Автоматизация деятельности гидов, инструкторов и туристических организаций по созданию, визуализации, публикации и сопровождению туристических маршрутов с мультимедийным контентом.

\vspace{4pt}

\textbf{Основные функции}

\begin{itemize}
\item регистрация и авторизация пользователей с ролевой моделью;
\item создание маршрутов на интерактивной карте и прикрепление геопривязанных фотографий;
\item автоматическое построение маршрута, расчёт расстояния, набора высоты, времени прохождения и сложности;
\item асинхронная обработка фотографий и хранение медиафайлов;
\item публикация маршрутов, поиск, фильтрация, социальные функции и публичный просмотр по ссылке;
\item ИИ-ассистент для геокодирования, поиска мест, погодных и маршрутных подсказок;
\item административная панель и средства модерации.
\end{itemize}

\textbf{Используемые технологии и платформы}

Rust, Go, TypeScript, React, Tauri, PostgreSQL, NATS JetStream, Redis, MinIO, Docker, Kubernetes, GitHub Actions, OpenStreetMap/Nominatim, OSRM, API языковых моделей.

\vspace{8pt}

\textbf{ВЫПОЛНЕНИЕ РАБОТЫ}

\vspace{4pt}

\textbf{Решаемые задачи}

\begin{enumerate}
\item проанализировать предметную область цифрового туризма и существующие аналоги;
\item сформулировать функциональные и нефункциональные требования к системе;
\item спроектировать архитектуру системы, схему базы данных и пользовательские интерфейсы;
\item реализовать серверную часть системы в виде набора взаимодействующих сервисов;
\item разработать клиентское веб-приложение и настольную версию;
\item реализовать обработку геопривязанных фотографий и интеграцию ИИ-функций;
\item настроить контейнеризацию, CI/CD и развёртывание;
\item провести тестирование, оценить производительность и подготовить документацию.
\end{enumerate}

\textbf{Состав технической документации}

\begin{itemize}
\item пояснительная записка к ВКР;
\item техническое задание;
\item программа и методика испытаний;
\item руководство пользователя и администратора в составе пояснительной записки.
\end{itemize}

\textbf{Состав графической части}

\begin{itemize}
\item диаграмма бизнес-процесса;
\item диаграмма контекста системы;
\item диаграмма контейнеров;
\item ER-диаграмма базы данных;
\item макеты ключевых экранов и материалы презентации к защите.
\end{itemize}

\vspace{8pt}

\textbf{ПЛАН РАБОТЫ НАД ВКР}

\vspace{6pt}

\begin{landscape}
\small
\begin{center}
\begin{tabular}{|p{7.2cm}|*{18}{>{\centering\arraybackslash}p{0.6cm}|}}
\hline
\textbf{ЗАДАЧИ} & \multicolumn{18}{c|}{\textbf{НЕДЕЛИ}} \\
\cline{2-19}
& \textbf{1} & \textbf{2} & \textbf{3} & \textbf{4} & \textbf{5} & \textbf{6} & \textbf{7} & \textbf{8} & \textbf{9} & \textbf{10} & \textbf{11} & \textbf{12} & \textbf{13} & \textbf{14} & \textbf{15} & \textbf{16} & \textbf{17} & \textbf{18} \\
\hline
Анализ предметной области и аналогов
& \wk & \wk & \blankwk & \blankwk & \blankwk & \blankwk & \blankwk & \blankwk & \blankwk & \blankwk & \blankwk & \blankwk & \blankwk & \blankwk & \blankwk & \blankwk & \blankwk & \blankwk \\
\hline
Формирование требований к системе
& \blankwk & \wk & \wk & \blankwk & \blankwk & \blankwk & \blankwk & \blankwk & \blankwk & \blankwk & \blankwk & \blankwk & \blankwk & \blankwk & \blankwk & \blankwk & \blankwk & \blankwk \\
\hline
Проектирование архитектуры, БД и интерфейсов
& \blankwk & \blankwk & \wk & \wk & \wk & \blankwk & \blankwk & \blankwk & \blankwk & \blankwk & \blankwk & \blankwk & \blankwk & \blankwk & \blankwk & \blankwk & \blankwk & \blankwk \\
\hline
Реализация сервисов аутентификации и маршрутов
& \blankwk & \blankwk & \blankwk & \blankwk & \wk & \wk & \wk & \wk & \blankwk & \blankwk & \blankwk & \blankwk & \blankwk & \blankwk & \blankwk & \blankwk & \blankwk & \blankwk \\
\hline
Реализация обработки фотографий и ИИ-функций
& \blankwk & \blankwk & \blankwk & \blankwk & \blankwk & \blankwk & \blankwk & \wk & \wk & \wk & \blankwk & \blankwk & \blankwk & \blankwk & \blankwk & \blankwk & \blankwk & \blankwk \\
\hline
Разработка клиентского веб-приложения и desktop-версии
& \blankwk & \blankwk & \blankwk & \blankwk & \blankwk & \blankwk & \blankwk & \blankwk & \wk & \wk & \wk & \wk & \blankwk & \blankwk & \blankwk & \blankwk & \blankwk & \blankwk \\
\hline
Настройка Docker, Kubernetes и CI/CD
& \blankwk & \blankwk & \blankwk & \blankwk & \blankwk & \blankwk & \blankwk & \blankwk & \blankwk & \blankwk & \blankwk & \wk & \wk & \wk & \blankwk & \blankwk & \blankwk & \blankwk \\
\hline
Тестирование, устранение дефектов и анализ производительности
& \blankwk & \blankwk & \blankwk & \blankwk & \blankwk & \blankwk & \blankwk & \blankwk & \blankwk & \blankwk & \blankwk & \blankwk & \blankwk & \wk & \wk & \wk & \blankwk & \blankwk \\
\hline
Оформление пояснительной записки и материалов к защите
& \blankwk & \blankwk & \blankwk & \blankwk & \blankwk & \blankwk & \blankwk & \blankwk & \blankwk & \blankwk & \blankwk & \blankwk & \blankwk & \blankwk & \wk & \wk & \wk & \wk \\
\hline
\end{tabular}
\end{center}
\end{landscape}

\vspace{8pt}

\textbf{РУКОВОДИТЕЛЬ ОП:}\\
«\_\_\_» \_\_\_\_\_\_\_\_\_\_ 2026, \_\_\_\_\_\_\_\_\_\_\_\_\_\_\_ / А.\,Ю.~Филиппович, к.э.н. /\\
подпись \hspace{3.2cm} ФИО, уч. звание и степень

\vspace{10pt}

\textbf{РУКОВОДИТЕЛЬ ВКР:}\\
«\_\_\_» \_\_\_\_\_\_\_\_\_\_ 2026, \_\_\_\_\_\_\_\_\_\_\_\_\_\_\_ / Васильев~Д.\,Б., старший преподаватель /\\
подпись \hspace{3.2cm} ФИО, уч. звание и степень

\vspace{10pt}

\textbf{СТУДЕНТ:}\\
«\_\_\_» \_\_\_\_\_\_\_\_\_\_ 2026, \_\_\_\_\_\_\_\_\_\_\_\_\_\_\_ / Дубровских~Н.\,Е., 221-361 /\\
подпись \hspace{3.2cm} ФИО, группа
